% !TeX root = ../Report.tex

\chapter{Introduction}

Space mission operations require the continuous coordination and interpretation of large amounts of heterogeneous information. Within a modern ground segment, operational data are distributed across multiple systems responsible for mission planning, spacecraft monitoring, telemetry, network status, anomaly and incident management, and procedural knowledge. Although these systems are individually specialized, mission operations emerge from their continuous interaction. Consequently, maintaining an accurate and coherent view of the overall mission state requires operators to correlate information across multiple sources, often under time constraints and in the presence of incomplete or conflicting evidence. This places a substantial cognitive burden on mission operators, particularly during unexpected events and other situations requiring rapid operational assessment.

The Mission Manager is particularly exposed to this information complexity. The role requires maintaining situational awareness, evaluating operational events, coordinating responses to anomalies, and supporting decisions during both nominal and contingency operations. These activities involve integrating information originating from different operational domains, including telemetry, mission planning, flight dynamics, communication systems, procedures, and operator reports. The increasing volume and heterogeneity of such information motivate the adoption of intelligent Decision Support Systems (DSS) capable of aggregating relevant evidence, reducing information fragmentation, and assisting operators in interpreting the current operational context.

Recent advances in Large Language Models (LLMs) have created new opportunities for AI-based decision support. Their ability to process natural language, interact with external tools, and integrate information from multiple sources makes them promising components for complex operational environments. However, directly providing a general-purpose LLM with a large operational context does not necessarily ensure that the model will identify and use the most relevant information effectively. In particular, LLM-based systems remain affected by hallucinations, irrelevant context, and limitations in evidence aggregation, while the complexity of mission operations requires responses to remain grounded in the available operational information. The literature reviewed in this work also shows that existing LLM-based decision-support systems for space operations are largely domain-specific and prototype-oriented, often addressing individual scenarios rather than providing a general multi-agent architecture capable of coordinating heterogeneous operational information and reasoning processes.

This thesis addresses this problem through the development of FUNES (Foresight Understanding for Nominal and Emergency Scenarios), an AI-based Decision Support System designed for space mission operations. FUNES acts as an intermediate layer between existing mission operational systems and human operators, collecting heterogeneous operational data and providing context-aware decision support. Rather than relying on a single general-purpose LLM to perform the entire investigation process, the proposed architecture decomposes the task into specialized agents responsible for planning information acquisition, querying operational sources, filtering retrieved evidence, synthesizing the final assessment, and evaluating whether additional investigation is required. These agents are coordinated through a graph-based orchestration mechanism and operate according to an evidence-driven approach.
The main objective of the work is therefore not to replace the reasoning capability of the underlying language model, but to improve how heterogeneous operational information is selected, organized, and provided to it. The proposed architecture introduces explicit mechanisms for information planning, source-specific retrieval, evidence refinement, iterative investigation, and response synthesis. In this sense, FUNES is intended to act as a control and orchestration layer around the LLM, reducing irrelevant contextual information while preserving the evidence required to support an operational assessment.

To investigate the effectiveness of this approach, a Proof of Concept was developed and evaluated against a general-purpose Llama 3.3 70B Instruct model operating directly on the complete operational context. The evaluation consists of twenty operational questions organized into three levels of complexity: single-source information retrieval, multi-source information correlation, and decision-support tasks requiring anomaly identification, impact assessment, and operational recommendations. The results show that FUNES achieved an overall score of 89.9\%, compared with 63.0\% for the baseline, corresponding to an improvement of 26.9 percentage points. The proposed system achieved higher aggregate scores across all three task-complexity categories, with the largest improvement observed for retrieval-oriented tasks.
At the same time, the evaluation shows that improved information retrieval and evidence correlation do not automatically guarantee superior operational reasoning. In some complex tasks, FUNES retrieved the relevant evidence but did not fully translate it into the specific operational information requested by the user. These results highlight the importance of the final synthesis stage and indicate that the proposed architecture should be understood primarily as a mechanism for controlling and structuring the information available to the underlying LLM, rather than as a replacement for its reasoning capabilities.

The remainder of this thesis is organized as follows. Chapter 2 reviews the evolution of Natural Language Processing and Decision Support Systems, with particular emphasis on Transformers, Large Language Models, hallucination and reliability issues, Retrieval-Augmented Generation, reasoning frameworks, and agentic architectures. Chapter 3 presents the system engineering and requirement analysis of FUNES, including the mission context, system architecture, operational use cases, and model selection. Chapter 4 describes the development of the Proof of Concept. Chapter 5 presents the experimental evaluation, results, discussion, and limitations of the proposed system. Chapter 6 discusses future research directions, while Chapter 7 concludes the thesis.

